% WRITE LAST
\section{Introduction}
\label{sec:intro}
Skip Connections-NN (SCNN) wanted to solve grad vanishing problem that VGG introduced. By defining $x = f(x) + x$ it induced some properties: Works in the same feature representation. P.e. A given channel
will have the same purpose in the next layer.

\citep{sander2022residualneuralnetworksdiscretize, queiruga2020continuousindepthneuralnetworks} says that resnets just have the algebraic form of Euler ODE solver, that if params are steps they won'r be ODEs.
We apport evidence that residual blocks are 1-step ODE solvers and we therefore have a sum of Ode solvers that are constantly changing of domain and we can accomplish the same dynamics by sharing the weights and
we augmentate the resolution of the feature dynamics to study with better quality how X moves along the latent-space manifold in time / depth.

Questions we want to solve:

\begin{itemize}
    \item Can SCNN's blocks be looped indefinitely?
        $\rightarrow$ Yes
    \item Are all this weights necessary? + Shared networks can reach similar performance?
        $\rightarrow$ No + Yes
    \item What are the dynamics of the features in this two types of models?
    Statement: The over-parametrized networks undo what the previous layer does, that's why SCNNs and weight-tied are equal.
        Tangential acceleration is crucial for this affirmation.
        Here we analyze the geometry of trajectories.
        Take into account the massive activations here.
    \item How can we improve discrete models given this evidence? -> Delta networks?
\end{itemize}
